%==============================================================================
% preamble.tex
% Shared packages, styling and math macros for the Compendium.
% Chapters should never load packages themselves — add them here.
%==============================================================================

%------------------------------------------------------------------ encoding --
\usepackage[T1]{fontenc}
\usepackage[utf8]{inputenc}
\usepackage[english]{babel}
\usepackage{lmodern}
\usepackage{microtype}

%------------------------------------------------------------------- layout ---
\usepackage[a4paper,margin=2.5cm,headheight=14pt]{geometry}
\usepackage{fancyhdr}
\usepackage{titlesec}

%--------------------------------------------------------------------- math ---
\usepackage{amsmath}
\usepackage{amssymb}
\usepackage{amsthm}
\usepackage{mathtools}
\usepackage{bm}

%------------------------------------------------------------------ graphics --
\usepackage{graphicx}
\usepackage{xcolor}
\usepackage{float}
\usepackage[font=small,labelfont=bf]{caption}
\usepackage{subcaption}
\graphicspath{{figures/}}

%------------------------------------------------------------------- content --
\usepackage{booktabs}
\usepackage{enumitem}
\usepackage{listings}
\usepackage[most]{tcolorbox}

%--------------------------------------------------------------- references ---
\usepackage{csquotes}   % required by biblatex when babel is loaded
\usepackage[backend=biber,style=numeric,sorting=nyt,maxbibnames=99]{biblatex}
\usepackage{hyperref}   % load late
\usepackage[capitalise,noabbrev]{cleveref}  % load after hyperref

%==============================================================================
% Colour palette
%==============================================================================
\definecolor{cmpNavy}{HTML}{1F3A5F}
\definecolor{cmpTeal}{HTML}{18827E}
\definecolor{cmpAmber}{HTML}{B26B00}
\definecolor{cmpGrey}{HTML}{4A4A4A}
\definecolor{cmpMist}{HTML}{F4F6F8}

\hypersetup{
    colorlinks = true,
    linkcolor  = cmpNavy,
    citecolor  = cmpTeal,
    urlcolor   = cmpTeal,
    pdftitle   = {The Applied AI Engineer's Compendium},
    pdfauthor  = {Euan Cortes},
    bookmarksnumbered = true,
}

%==============================================================================
% Headers, footers and chapter titles
%==============================================================================
\pagestyle{fancy}
\fancyhf{}
\fancyhead[L]{\nouppercase{\small\textsc{\leftmark}}}
\fancyhead[R]{\small\thepage}
\renewcommand{\headrulewidth}{0.4pt}

\fancypagestyle{plain}{%
    \fancyhf{}%
    \fancyfoot[C]{\small\thepage}%
    \renewcommand{\headrulewidth}{0pt}%
}

\titleformat{\chapter}[display]
    {\normalfont\huge\bfseries\color{cmpNavy}}
    {\filleft\Large\mdseries\scshape Chapter \thechapter}
    {2ex}
    {\titlerule\vspace{2ex}\filright}
    [\vspace{1ex}{\titlerule[0.6pt]}]

\titleformat{\section}
    {\normalfont\Large\bfseries\color{cmpNavy}}{\thesection}{1em}{}
\titleformat{\subsection}
    {\normalfont\large\bfseries\color{cmpGrey}}{\thesubsection}{1em}{}

\setcounter{secnumdepth}{3}
\setcounter{tocdepth}{2}

%==============================================================================
% Theorem-like environments
%
%   definition / theorem / lemma / proposition  — numbered within the chapter
%   example / remark                            — unnumbered, informal
%   intuition                                   — the "why does this matter" box
%==============================================================================
\tcbuselibrary{theorems,skins,breakable}

\newtcbtheorem[number within=chapter]{definition}{Definition}{%
    enhanced, breakable, colback=cmpMist, colframe=cmpNavy,
    fonttitle=\bfseries, boxrule=0.6pt, arc=2pt,
    attach boxed title to top left={yshift=-2mm, xshift=4mm},
    boxed title style={colback=cmpNavy},
}{def}

\newtcbtheorem[number within=chapter]{theorem}{Theorem}{%
    enhanced, breakable, colback=white, colframe=cmpTeal,
    fonttitle=\bfseries, boxrule=0.8pt, arc=2pt,
    attach boxed title to top left={yshift=-2mm, xshift=4mm},
    boxed title style={colback=cmpTeal},
}{thm}

\newtcbtheorem[use counter from=theorem]{lemma}{Lemma}{%
    enhanced, breakable, colback=white, colframe=cmpTeal,
    fonttitle=\bfseries, boxrule=0.8pt, arc=2pt,
    attach boxed title to top left={yshift=-2mm, xshift=4mm},
    boxed title style={colback=cmpTeal},
}{lem}

\newtcolorbox{intuition}[1][Intuition]{%
    enhanced, breakable, colback=cmpAmber!6, colframe=cmpAmber,
    fonttitle=\bfseries, boxrule=0.6pt, arc=2pt, left=3mm, right=3mm,
    title={#1},
}

\theoremstyle{definition}
\newtheorem{example}{Example}[chapter]
\theoremstyle{remark}
\newtheorem*{remark}{Remark}

% cleveref names for the tcolorbox theorems
\crefname{def}{Definition}{Definitions}
\crefname{thm}{Theorem}{Theorems}
\crefname{lem}{Lemma}{Lemmas}

%==============================================================================
% Code listings (Python — the applied half of the Compendium)
%==============================================================================
\definecolor{lstKeyword}{HTML}{1F3A5F}
\definecolor{lstString}{HTML}{18827E}
\definecolor{lstComment}{HTML}{8A8A8A}

\lstdefinestyle{python}{
    language         = Python,
    basicstyle       = \ttfamily\small,
    keywordstyle     = \color{lstKeyword}\bfseries,
    stringstyle      = \color{lstString},
    commentstyle     = \color{lstComment}\itshape,
    numbers          = left,
    numberstyle      = \tiny\color{cmpGrey},
    numbersep        = 8pt,
    frame            = single,
    rulecolor        = \color{cmpGrey!40},
    backgroundcolor  = \color{cmpMist},
    breaklines       = true,
    showstringspaces = false,
    tabsize          = 4,
    captionpos       = b,
}
\lstset{style=python}
\renewcommand{\lstlistingname}{Listing}

%==============================================================================
% Math macros
% Keep every notational choice here so notation stays consistent book-wide.
%==============================================================================
% Sets and spaces
\newcommand{\R}{\mathbb{R}}
\newcommand{\N}{\mathbb{N}}
\newcommand{\Z}{\mathbb{Z}}
\newcommand{\Prob}{\mathbb{P}}

% Vectors, matrices, tensors
\newcommand{\vect}[1]{\bm{#1}}
\newcommand{\mat}[1]{\bm{\mathrm{#1}}}
\newcommand{\T}{^{\top}}

% Operators
\DeclareMathOperator*{\argmin}{arg\,min}
\DeclareMathOperator*{\argmax}{arg\,max}
\DeclareMathOperator{\diag}{diag}
\DeclareMathOperator{\rank}{rank}
\DeclareMathOperator{\tr}{tr}
\DeclareMathOperator{\Var}{Var}
\DeclareMathOperator{\Cov}{Cov}
\DeclareMathOperator{\softmax}{softmax}
\DeclareMathOperator{\sgn}{sgn}
\newcommand{\E}{\mathbb{E}}

% Delimiters
\DeclarePairedDelimiter{\abs}{\lvert}{\rvert}
\DeclarePairedDelimiter{\norm}{\lVert}{\rVert}
\DeclarePairedDelimiter{\inner}{\langle}{\rangle}
\newcommand{\given}{\mathbin{\vert}}

% Calculus
\newcommand{\grad}{\nabla}
\newcommand{\pd}[2]{\frac{\partial #1}{\partial #2}}
\newcommand{\dd}[2]{\frac{\mathrm{d} #1}{\mathrm{d} #2}}
\newcommand{\Hess}{\mat{H}}

% Distributions and divergences
\newcommand{\Normal}{\mathcal{N}}
\newcommand{\KL}[2]{D_{\mathrm{KL}}\!\left(#1 \,\|\, #2\right)}
\newcommand{\Ent}[1]{H\!\left(#1\right)}
\newcommand{\dist}{\sim}

% Cross-referencing helper: link a theory section to its project folder
\newcommand{\projectref}[1]{%
    \marginpar{\footnotesize\raggedright\textcolor{cmpTeal}{$\blacktriangleright$}\;\texttt{#1}}%
}
